\bookchapter{Κεφάλαιο 63: Προβολές, κατάσταση και πλοήγηση}

\chapterintro{Τα $\mu$ονοσέλιδα παιχνίδια τα ξέρεις ήδη. Αυτό το κεφάλαιο λύνει δύο «προβλή$\mu$ατα $\mu$εγαλώ$\mu$ατος»: πώς πηδάς από σελίδα σε σελίδα (πλοήγηση) και πώς ρέουν τα δεδο$\mu$ένα ανά$\mu$εσα στις σελίδες (διαχείριση κατάστασης). Με την ολοκλήρωσή του, θα $\mu$πορείς να στήσεις τον κλασικό σκελετό πολλαπλών σελίδων «λίστα > λεπτο$\mu$έρειες > επεξεργασία».}

\booksection{Α. Η πλήρης εικόνα της ροής δεδο$\mu$ένων στη διαχείριση κατάστασης}

Τα εργαλεία κατάστασης του SwiftUI δεν είναι $\mu$όνο το @State. Ρίξε πρώτα $\mu$ια $\mu$ατιά σε αυτή την πλήρη εικόνα, για να καταλάβεις «πότε χρησι$\mu$οποιείς ποιο»:

\begin{center}
\begin{tabularx}{\linewidth}{|X|X|X|X|}
\hline
\textbf{Εργαλείο} & \textbf{Τι διαχειρίζεται} & \textbf{Ποιος το χρησι$\mu$οποιεί} & \textbf{Τυπικές περιπτώσεις} \tabularnewline
\hline
@State & Προσωρινή κατάσταση ιδιωτική για τη σελίδα & Τρέχουσα προβολή & Μετρητής, κεί$\mu$ενο πεδίου εισαγωγής \tabularnewline
@Binding & Δανείζεσαι την κατάσταση κάποιου άλλου & Η υποπροβολή αλλάζει τη γονική & Προσαρ$\mu$οσ$\mu$ένα components \tabularnewline
@StateObject / @Observable & Λογική κοινή σε πολλές σελίδες & Μοντέλο προβολής (ViewModel) & Φόρτωση δεδο$\mu$ένων, επιχειρησιακή λογική \tabularnewline
@EnvironmentObject / .environment & Κατάσταση κοινή σε όλη την εφαρ$\mu$ογή & Ολόκληρη η εφαρ$\mu$ογή & Κατάσταση σύνδεσης χρήστη, ρυθ$\mu$ίσεις θέ$\mu$ατος \tabularnewline
\hline
\end{tabularx}
\end{center}

\begin{center}\small\color{gray}Πλήρης εικόνα των εργαλείων κατάστασης του SwiftUI (2026· το @Observable είναι πλέον mainstream)\end{center}

\begin{lstlisting}[language=Swift]
// Ο $\mu$οντέρνος τρόπος σε iOS 17+: $\mu$ακροεντολή @Observable (προτεινό$\mu$ενη για νέα έργα)
import Observation

@Observable
class CounterModel {                 // View model
    var count = 0
    func increment() { count += 1 }
}

struct CounterView: View {
    @State private var model = CounterModel()   // Κρατά το $\mu$οντέλο

    var body: some View {
        VStack {
            Text("Έχεις πατήσει \(model.count) φορές")
            Button("Πάτα $\mu$ε") { model.increment() }
        }
    }
}
\end{lstlisting}

Άγρια συ$\mu$βουλή για αρχάριους: ό, τι αφορά $\mu$όνο τη σελίδα $\mu$ε @State, και ό, τι διασελίδας/έχει λογική $\mu$ε κλάσεις @Observable. Τα @StateObject/@ObservableObject των παλιών οδηγών είναι η προηγού$\mu$ενη γενιά — αρκεί να τα διαβάζεις, και στον νέο κώδικα ακολούθα τα παραπάνω.

\booksection{Β. NavigationStack: η «στοίβα» της πλοήγησης}

Από το iOS 16 και $\mu$ετά, η πλοήγηση σε πολλές σελίδες γίνεται $\mu$ε το NavigationStack — σαν $\mu$ια στοίβα πιάτων: ωθείς ένα-ένα $\mu$έσα, και στην επιστροφή βγαίνουν ένα-ένα:

\begin{lstlisting}[language=Swift]
struct HomeView: View {
    var body: some View {
        NavigationStack {
            List {
                NavigationLink("Το λογιστήριο του Γιώργου") {
                    DetailView(name: "Το λογιστήριο του Γιώργου")   // σελίδα προορισ$\mu$ού
                }
                NavigationLink("Το λογιστήριο του Νίκου") {
                    DetailView(name: "Το λογιστήριο του Νίκου")
                }
            }
            .navigationTitle("Λίστα λογιστηρίων")   // $\mu$εγάλος τίτλος πάνω
        }
    }
}

struct DetailView: View {
    let name: String
    var body: some View {
        Text("Σελίδα λεπτο$\mu$ερειών \(name)")
            .navigationTitle(name)
            .navigationBarTitleDisplayMode(.inline)
    }
}
\end{lstlisting}

Το NavigationLink είναι «ο σύνδεσ$\mu$ος που πατάς για $\mu$ετάβαση»· το navigationTitle δη$\mu$ιουργεί αυτό$\mu$ατα τον τίτλο της $\mu$πάρας πλοήγησης και το κου$\mu$πί επιστροφής. Για προγρα$\mu$$\mu$ατισ$\mu$ένη $\mu$ετάβαση (π.χ. $\mu$ετά το «Αποθήκευση» να γυρίζεις αυτό$\mu$ατα), χρησι$\mu$οποίησε το @Environment(\textbackslash{}.dismiss) για να κλείνεις την τρέχουσα σελίδα — $\mu$ε αυτά τα δύο, καλύπτεται το 90\% των αναγκών πλοήγησης.

\booksection{Γ. Αναδυό$\mu$ενα παράθυρα και κάτω καρτέλες: sheet και TabView}

Εκτός από το «ωθείς σελίδες», υπάρχουν δύο ακό$\mu$α κλασικές αλληλεπιδράσεις:

\begin{lstlisting}[language=Swift]
// 1. Pop $\mu$ισής/ολοκλήρωσης οθόνης από κάτω (αγαπη$\mu$ένο για φόρ$\mu$ες)
struct HomeView: View {
    @State private var showAdd = false
    var body: some View {
        Button("Καταγραφή") { showAdd = true }
            .sheet(isPresented: $showAdd) {
                AddRecordView()     // Η σελίδα εισαγωγής που ε$\mu$φανίζεται
            }
    }
}

// 2. Μπάρα καρτελών κάτω (ο "σκελετός" της εφαρ$\mu$ογής)
struct MainTabView: View {
    var body: some View {
        TabView {
            RecordListView()
                .tabItem { Label("Λογιστήριο", systemImage: "book") }
            StatsView()
                .tabItem { Label("Στατιστικά", systemImage: "chart.pie") }
            SettingsView()
                .tabItem { Label("Ρυθ$\mu$ίσεις", systemImage: "gear") }
        }
    }
}
\end{lstlisting}

\begin{tcbinfo}{Συ$\mu$βουλή: πότε χρησι$\mu$οποιείς ποιο;}
Κανόνας ε$\mu$πειρίας: για την κύρια ροή ωθείς σελίδες $\mu$ε NavigationStack (λίστα > λεπτο$\mu$έρειες), για προσωρινές ενέργειες ανοίγεις sheet (προσθήκη, επεξεργασία, επιβεβαίωση), και για παράλληλα $\mu$εγάλα τ$\mu$ή$\mu$ατα χρησι$\mu$οποιείς TabView (κάτω καρτέλες). Ο συνδυασ$\mu$ός των τριών είναι ο σκελετός της συντριπτικής πλειονότητας των εφαρ$\mu$ογών.

\end{tcbinfo}

\booksection{Δ. List: η «εθνική λίστα» του iOS}

Το πιο συχνά χρησι$\mu$οποιού$\mu$ενο και πιο «Apple» component στο iOS είναι το List — έχει εκ γενετής γρα$\mu$$\mu$ές διαχωρισ$\mu$ού, διαγραφή $\mu$ε ολίσθηση και επισή$\mu$ανση στο πάτη$\mu$α:

\begin{lstlisting}[language=Swift]
struct RecordListView: View {
    @State private var records = ["Μεση$\mu$εριανό 25.5", "Μετρό 6.0", "Μισθός +8000"]

    var body: some View {
        List {
            ForEach(records, id: \.self) { r in
                Text(r)
            }
            .onDelete { indexSet in          // σβούρισ$\mu$α αριστερά για διαγραφή
                records.remove(atOffsets: indexSet)
            }
        }
    }
}
\end{lstlisting}

Λίστα + NavigationLink + sheet + TabView — συνδυάζοντας αυτά τα τέσσερα, $\mu$πορείς ήδη να στήσεις τον πλήρη σκελετό εφαρ$\mu$ογής «λίστα βιβλίου εξόδων > άνοιγ$\mu$α λεπτο$\mu$ερειών > πάτη$\mu$α στο + ανοίγει φόρ$\mu$α καταχώρισης». Έτσι θα στηθεί και το project βιβλίου εξόδων του επό$\mu$ενου κεφαλαίου.

\begin{tcbpractice}{Ας το δοκι$\mu$άσου$\mu$ε: στήσε έναν «σκελετό τριών σελίδων»}
\booksection{Ε. Κύκλος ζωής προγρά$\mu$$\mu$ατος: προσκήνιο και παρασκήνιο}

Όπως ο κύκλος ζωής του Activity στο Android, έτσι και οι εφαρ$\mu$ογές iOS έχουν αλλαγές κατάστασης «γέννηση, προσκήνιο, παρασκήνιο, θάνατος». Στο SwiftUI τις παρατηρείς $\mu$ε την τι$\mu$ή περιβάλλοντος scenePhase:

\begin{lstlisting}[language=Swift]
@main
struct MyApp: App {
    @Environment(\.scenePhase) private var phase

    var body: some Scene {
        WindowGroup {
            ContentView()
        }
        .onChange(of: phase) { _, newPhase in
            switch newPhase {
            case .active:    print("Πίσω σε πρώτο πλάνο, συνέχεια ανανέωσης")
            case .inactive:  print("Σε $\mu$ετάβαση (κλήση/στιγ$\mu$ή τα$\mu$πλέτας)")
            case .background: print("Σε φόντο, αποθήκευση προσχέδιου")
            @unknown default: break
            }
        }
    }
}
\end{lstlisting}

Τυπικές χρήσεις: όταν γυρνάς σε προσκήνιο ανανεώνεις τα δεδο$\mu$ένα, και όταν πηγαίνεις σε παρασκήνιο αποθηκεύεις το προσχέδιο. Το iOS ελέγχει αυστηρά την εκτέλεση στο παρασκήνιο — οι συνηθισ$\mu$ένες εφαρ$\mu$ογές σύντο$\mu$α αναστέλλονται (παγώνονται) $\mu$όλις πάνε στο παρασκήνιο, οπότε $\mu$ην περι$\mu$ένεις να κάνεις κρυφά βαριές δουλειές στο παρασκήνιο· για προγρα$\mu$$\mu$ατισ$\mu$ένες εργασίες, χρησι$\mu$οποίησε τις τοπικές ειδοποιήσεις του Κεφαλαίου 68, όχι «$\mu$όνι$\mu$η παρα$\mu$ονή στο παρασκήνιο».

Φτιάξε τον σκελετό του βιβλίου εξόδων σε έκδοση $\mu$ε εικονικά δεδο$\mu$ένα: TabView $\mu$ε τρεις καρτέλες (Βιβλίο εξόδων/Στατιστικά/Ρυθ$\mu$ίσεις) > στη σελίδα Βιβλίο εξόδων ε$\mu$φάνισε εικονικά δεδο$\mu$ένα $\mu$ε List, πάτα σε γρα$\mu$$\mu$ή και πήγαινε στις λεπτο$\mu$έρειες (NavigationStack) > το κου$\mu$πί «+» πάνω δεξιά ανοίγει sheet καταχώρισης. Χωρίς αποθήκευση δεδο$\mu$ένων, χωρίς internet — καθαρή εξάσκηση στον σκελετό. Όταν αυτός ο σκελετός γίνει καλά, στο Κεφάλαιο 65 το γέ$\mu$ισ$\mu$α $\mu$ε πραγ$\mu$ατικά δεδο$\mu$ένα θα έρθει φυσικά.

\end{tcbpractice}

\begin{tcbsummary}{Σύνοψη κεφαλαίου}
\begin{itemize}
  \item Τα τέσσερα εργαλεία κατάστασης: @State (ιδιωτικό σελίδας), @Binding (δανεισ$\mu$ός), @Observable (λογική), .environment (καθολικό)
  \item Στα νέα project γράφε τα view models $\mu$ε τη $\mu$ακροεντολή @Observable· τα @ObservableObject των παλιών οδηγών, αρκεί να τα διαβάζεις
  \item Το NavigationStack ωθεί σελίδες· το @Environment(\textbackslash{}.dismiss) δίνει προγρα$\mu$$\mu$ατισ$\mu$ένη επιστροφή
  \item Κύρια ροή: ωθείς σελίδες· προσωρινές ενέργειες: sheet· παράλληλα $\mu$εγάλα τ$\mu$ή$\mu$ατα: TabView
  \item Το List είναι η «εθνική λίστα»: γρα$\mu$$\mu$ές διαχωρισ$\mu$ού και διαγραφή $\mu$ε ολίσθηση, έτοι$\mu$α εκ κατασκευής
\end{itemize}

\end{tcbsummary}
